\section{Simulation in the DCP representation}\label{sec:simulator-dcp-representation}
In this section, we describe how to update the DCP-stabilizer group of a state under the operations of preparing $|0\rangle$ and $|+\rangle$ states, applying Pauli, $\CNOT$ and $\T$ gates, and measuring $Z$ and certain $X$ Pauli operators projectively. We show that each of these operations can be simulated in polynomial time using the canonical representation of the DCP-stabilizer group.

First, note that any product state $\ket{0}^{\otimes a} \otimes \ket{+}^{\otimes b}$ is a DCP state, since it is stabilized by the DCP group generated by $\langle Z_1,\dots,Z_a, X_{a+1},\dots,X_{a+b} \rangle$. Thus, we can initialize the DCP-stabilizer group of any product state in polynomial time. The other operations are more complicated and described in the following subsections.

\subsection{Unitary operations}
Let's consider the following very general fact about describing states in the Heisenberg picture, that is, by considering a set of matrices that stabilize the state:
\begin{proposition}\label{prop:heisenberg}
    Suppose that we have a state $\ket{\psi}$ and a set of matrices $S$ that stabilize it. Consider a unitary operation $U$ that is applied to the state. Then, the set of matrices $U S U^\dagger$ stabilizes the state $U\ket{\psi}$, where
    \[
        U S U^\dagger = \{ U g U^\dagger \mid g \in S \}.
    \]
    Moreover, $\dim \Fix(S) = \dim \Fix(U S U^\dagger)$.
\end{proposition}
\begin{proof}
    We have that for any $L \in S$,
    \[
        U L U^\dagger (U \ket{\psi}) = U L \ket{\psi} = U \ket{\psi},
    \]
    so $U S U^\dagger$ stabilizes $U\ket{\psi}$.

    For the second part, note that if $\{v^{(1)}, \dots, v^{(r)}\}$ is a basis for $\Fix(S)$, then $\{U v^{(1)}, \dots, U v^{(r)}\}$ remains linearly independent and lies in $\Fix(U S U^\dagger)$, so $\dim \Fix(S) \le \dim \Fix(U S U^\dagger)$. The same argument applies to $U^\dagger$, so we also have $\dim \Fix(U S U^\dagger) \le \dim \Fix(S)$, and thus $\dim \Fix(S) = \dim \Fix(U S U^\dagger)$.
\end{proof}

By the conjugation-closure result \Cref{prop:dcp-closed}, we can apply \Cref{prop:heisenberg} to update the DCP-stabilizer group directly by conjugating its generators:

\begin{proposition}[Updating under Pauli, $\CNOT$, and $\T$ gates]\label{prop:dcp-update}
    Let $\ket{\psi}$ be a state with DCP-stabilizer group $\mathcal{D}(\psi) = \langle g_1, \dots, g_m \rangle$, and let $U$ be a Pauli operator, a $\CNOT$ gate, or a $\T$ (or $\T^\dagger$) gate. Then $Ug_1U^\dagger, \dots, Ug_mU^\dagger$ are DCP operators, by \Cref{prop:dcp-closed}, and
    \[
        \mathcal{D}(U\psi) \;=\; \big\langle\, Ug_1U^\dagger,\, \dots,\, Ug_mU^\dagger \,\big\rangle.
    \]
\end{proposition}
\begin{proof}
    Write $H = \langle Ug_1U^\dagger, \dots, Ug_mU^\dagger\rangle = U\,\mathcal{D}(\psi)\,U^\dagger$. For $h \in \mathcal{D}^n$, $h(U\psi) = U\psi$ iff $(U^\dagger h U)\psi = \psi$; since $U^\dagger$ is again a Pauli operator, $\CNOT$ gate, or $\T$-type gate, \Cref{prop:dcp-closed} guarantees $U^\dagger h U \in \mathcal{D}^n$ whenever $h \in \mathcal{D}^n$, so this is equivalent to $U^\dagger h U \in \mathcal{D}(\psi)$, i.e.\ to $h \in U\,\mathcal{D}(\psi)\,U^\dagger = H$. Hence $\mathcal{D}(U\psi) = H$.
\end{proof}

\subsection{Z-measurements}

When considering the issue of updating the DCP-stabilizer group under $Z$-measurements, the following is fundamental:
\begin{lemma}\label{lem:z-commute}
    Given a DCP operator $L$ and a product of Pauli $Z$ operators $z$, we have that $L$ and $z$ either commute or anticommute (i.e. $L z = \pm z L$).
\end{lemma}
\begin{proof}
    Write $L=\omega_8^kD(\vec{s},M)X(\vec{x})$ in the normal form of
    \Cref{prop:dcp-normal-form}, and write $z=Z(\vec{v})$ for some
    $\vec{v}\in\F_2^n$.  Since $D(\vec{s},M)$ and $z$ are both diagonal,
    they commute.  On the other hand, the Pauli commutation relations give
    \[
        X(\vec{x})Z(\vec{v})
        =(-1)^{\langle\vec{x},\vec{v}\rangle}
          Z(\vec{v})X(\vec{x}).
    \]
    It follows that
    $Lz=(-1)^{\langle\vec{x},\vec{v}\rangle}zL$, proving the claim.
\end{proof}

Just as in the Pauli stabilizer formalism, \Cref{lem:z-commute} lets us split the update rule for measuring a $Z$-string $z$ into a random and a deterministic case, according to how $z$ interacts with the current generators:

\newcommand{\dcpZMeasurementUpdateStatement}{%
    Let $\ket{\psi}$ be a DCP state, let
    $G=\langle g_1,\dots,g_m\rangle$ be a DCP-stabilizer group for
    $\ket{\psi}$, and let $z$ be a product of Pauli $Z$ operators.
    \begin{enumerate}[nosep,label=(\alph*)]
        \item \emph{(Random outcome.)} Suppose that $z$ anticommutes with
        some element of $G$, and let
        \[
            K=C_G(z)=\{g\in G:gz=zg\}.
        \]
        Measuring $z$ yields each outcome
        $\varepsilon\in\{+1,-1\}$ with probability $\tfrac12$, and
        \[
            H_\varepsilon=\langle K,\varepsilon z\rangle
        \]
        is a DCP-stabilizer group for the post-measurement state.  In the
        canonical representation of \Cref{def:canon-form}, a generating set
        for $K$ is obtained by retaining all diagonal generators, choosing
        an anticommuting lift $g_j$, and replacing every other anticommuting
        lift $g_i$ by $g_i g_j$.  Thus the post-measurement update amounts
        to replacing $g_j$ by $\varepsilon z$ after these row operations.
        \item \emph{(Deterministic outcome.)} If $z$ commutes with every
        element of $G$, the outcome $\varepsilon\in\{+1,-1\}$ is
        determined: it is the unique sign for which
        $\varepsilon z\in\mathcal{D}(\psi)$.  The state is unchanged, so
        $G$ remains a DCP-stabilizer group for it.
    \end{enumerate}
}

\begin{proposition}[Updating under $Z$-measurements]\label{prop:z-update}
\dcpZMeasurementUpdateStatement
\end{proposition}
The proof can be found in the \hyperref[appendix:tableau-update-measurements]{measurement update subsection} of \cref{appendix:dcp-tableaux}.

\subsection{X-measurements}
In the Heisenberg picture, the compatibility criteria and update rules for $X$-measurements closely mirror those of the phase-polynomial picture, described in \cref{thm:x-measurements}.
For a general $X(m)$ measurement, \cref{eq:x-measurement-reduction} lets us conjugate the tableau by $C_m$, perform the single-qubit $X_a$ update below, and conjugate back by $C_m^\dagger$. The identity measurement leaves the tableau unchanged.
Consider a 3PP state $\ket{\psi}$ with DCP-stabilizer group $G$, and suppose that we measure the single-qubit operator $X_a$, obtaining the outcome $\varepsilon = (-1)^\mu$ with $\mu \in \F_2$. Write $V = X(G)$, regarded as a vector space over $\F_2$.
Since $\ket{\psi}$ spans the one-dimensional space $\Fix(G)$, it is atomic for $G$ by \Cref{thm:atomic-basis}, so its support is a coset $S = x_0 \oplus V$ of $V$.
The two top-level cases of \cref{thm:x-measurements} are therefore exactly:
\begin{align}
    \text{(1)}\;\; e_a \not\in V
    \qquad\text{or}\qquad
    \text{(2)}\;\; e_a \in V.
\end{align}
Let us consider case (1) first. Measurements falling into this case are always compatible with the state and yield each outcome with probability $\tfrac12$, so all that remains is to describe how the tableau updates.

Suppose first that $v_a = 0$ for every $v \in V$.
Then qubit $a$ is unentangled from the rest of the system and sits in a computational basis state $\ket{\sigma}$, where $(-1)^\sigma$ is the deterministic outcome of measuring $Z_a$. The measurement simply replaces $\ket{\sigma}$ by $\ket{\pm}$, according to the sign of $\varepsilon$.
On the tableau, we first detach qubit $a$ from the generators, replacing each factor $\S_a$ by the scalar $i^\sigma$ and each factor $\CZ_{aj}$ by $Z_j^\sigma$, which does not change their action on the state, and then add the generator $\varepsilon X_a$.

The general case reduces to this one. Since $e_a \not\in V$, Gaussian elimination yields a vector $w \in \F_2^n$ with $w_a = 0$ and $\langle w, v \rangle = v_a$ for all $v \in V$. The circuit $C = \prod_{i \,:\, w_i = 1} \CNOT_{i,a}$ then clears the $a$-th coordinate of every element of $V$, and it commutes with $X_a$. Measuring $X_a$ is therefore equivalent to applying $C$, measuring $X_a$ as above, and then applying $C^\dagger$. Thus we conjugate the generators by $C$, update them as in the previous paragraph, and conjugate them back by $C^\dagger$, all of which takes polynomial time by \Cref{prop:dcp-closed}.

We now turn to case 2, where $e_a \in V$. By the surjectivity of $\pi$, there is some $g \in G$ of the form $g = \delta^* X(e_a)$, where $\delta^*$ is a diagonal DCP operator. Let $E : \F_2^n \to \Z_8$ denote the phase function of $\delta^*$, defined by $\delta^* \ket{x} = \omega_8^{E(x)} \ket{x}$. This function determines the discrete derivative on which \Cref{thm:x-measurements} is based:

\begin{proposition}\label{prop:derivative-from-E}
    Let $q$ be a phase polynomial for $\ket{\psi}$, so that, up to global phase,
    \[
        \ket{\psi} = \frac{1}{\sqrt{|S|}} \sum_{x \in S} \omega_8^{q(x)} \ket{x},
        \qquad S = x_0 \oplus V.
    \]
    Then, for every $x \in S$,
    \[
        \Delta_{e_a} q(x) = E(x \oplus e_a) = -E(x) \pmod 8.
    \]
\end{proposition}
\begin{proof}
    Since $e_a \in V$, we have $S \oplus e_a = S$. For $x \in S$,
    \[
        g\ket{x} = \delta^*\ket{x \oplus e_a} = \omega_8^{E(x \oplus e_a)}\ket{x \oplus e_a},
    \]
    so comparing the amplitudes of $\ket{x \oplus e_a}$ on both sides of $g\ket{\psi} = \ket{\psi}$ gives
    $q(x) + E(x \oplus e_a) = q(x \oplus e_a) \pmod 8$. Likewise, $g^{-1} = X(e_a)(\delta^*)^{-1}$ fixes $\ket{\psi}$ and
    \[
        g^{-1}\ket{x} = \omega_8^{-E(x)}\ket{x \oplus e_a},
    \]
    so that $q(x) - E(x) = q(x \oplus e_a) \pmod 8$.
\end{proof}

The derivative $\Delta_{e_a} q$ can thus be read off the DCP-stabilizer group, without tracking the phase polynomial itself. Substituting $\Delta_{e_a} q = -E$ into \Cref{thm:x-measurements} gives necessary and sufficient conditions for the compatibility of the measurement in terms of quantities computable from $G$, and \Cref{prop:x-measurement-case-decision} shows how to decide in polynomial time which of the cases 2a, 2b and 2c applies, if any. What remains is to describe how the DCP-stabilizer group is updated in each of them.

The starting point is the same in all three cases. Since $X_a = (\delta^*)^{-1} g$ and $g$ fix $\ket{\psi}$, the projection onto the $\varepsilon$-eigenspace of $X_a$ acts on $\ket{\psi}$ as a diagonal operator:
\begin{align}
    \frac{1}{2}(1 + \varepsilon X_a) \ket{\psi}
    &= \frac{1}{2}\big(1 + \varepsilon (\delta^*)^{-1} g\big) \ket{\psi} \notag \\
    &= \frac{1}{2}\big(1 + \varepsilon (\delta^*)^{-1}\big) \ket{\psi}.
    \label{eq:x-projection-diagonal}
\end{align}

In case 2a, $E(x) = -\Delta_{e_a} q(x) = 4f(x) \pmod 8$ on $S$ for an affine function $f$, which we write as $f(x) = c \oplus \langle u, x \rangle$ with $c \in \F_2$ and $u \in \F_2^n$.
Equivalently, $\delta^*$ acts on the support of $\ket{\psi}$ as the signed $Z$-string $(-1)^c Z(u)$.
This is the case, for instance, whenever $\delta^*$ is itself a signed product of $Z$ operators.
In general, $\delta^*$ may differ from $(-1)^c Z(u)$ by a diagonal operator that fixes every $\ket{x}$ with $x \in S$. As $(\delta^*)^{-1}\ket{\psi} = \delta^*\ket{\psi}$ in this case, \cref{eq:x-projection-diagonal} becomes
\begin{align}
    \frac{1}{2}(1 + \varepsilon X_a) \ket{\psi}
    &= \frac{1}{2}(1 + \varepsilon\, \delta^*) \ket{\psi} \\
    &= \frac{1}{2}\big(1 + (-1)^{\mu + c} Z(u)\big) \ket{\psi}.
\end{align}
Thus, measuring $X_a$ with outcome $\mu$ transforms the state exactly as measuring the $Z$-string $Z(u)$ with outcome $(-1)^{\mu + c}$. The outcome probabilities and the updated DCP-stabilizer group are therefore given by \Cref{prop:z-update}, and can be computed in polynomial time.

In cases 2b and 2c, the support is unchanged and, by \Cref{thm:x-measurements}, the difference $r = q' - q$ between the new and the old phase polynomial is itself a third-order phase polynomial. Writing it as a sum of terms,
\[
    r(x) = r_0 + \sum_i a_i x_i + 2\sum_{i<j} b_{ij} x_i x_j + 4\sum_{i<j<k} c_{ijk} x_i x_j x_k,
\]
adding $r$ to the phase polynomial has the same effect on the state, up to the global phase $\omega_8^{r_0}$, as applying the circuit
\[
    U =
    \prod_i \T_i^{a_i}
    \prod_{i<j} \CS_{ij}^{b_{ij}}
    \prod_{i<j<k} \CCZ_{ijk}^{c_{ijk}}.
\]
By \Cref{prop:heisenberg}, $U G U^\dagger$ is then a DCP-stabilizer group of the post-measurement state. Since $\CS$ and $\CCZ$ can be synthesized from $\CNOT$ and $\T$ gates, $U$ is a circuit of polynomially many such gates, so the new tableau can be computed in polynomial time by conjugating the generators gate by gate, as in \Cref{prop:dcp-update}.
